\chapter{الخدمات}
\label{ch:builder-services}

\section{ماذا ستبني في هذا الفصل؟}
\label{sec:builder-services-intro}

الخدمة هي الطريقة التي يصل بها نموذج أو أكثر إلى المتعاملين على بوابة المؤسسة، باسم واضح ودليل يقرؤونه قبل البدء، وشروط تحدد من يحق له التقديم، ورسوم ومسار معروف. ويشرح هذا الفصل إنشاء الخدمة من صفحة «الخدمات»، ثم ضبط تبويبات صفحتها واحداً بعد الآخر: التعريف، والتقديم والأهلية، والمسار والنماذج، والرسوم. ثم ينتقل إلى نشرها على البوابة، واستقبال الطلبات الحضورية في مراكز الخدمة، ومتابعة أداء الطلبات وأماكن تأخرها.

وتسري إعدادات الخدمة على البوابة فور حفظها دون إعادة نشر أي نموذج، أما التعديلات على النموذج نفسه فلا تسري إلا بعد نشره في منشئ النماذج. لذلك ابنِ النموذج وانشره أولاً، ثم أنشئ الخدمة حوله.

\begin{rcnote}[مثال: شركة الأفق]
تريد الأفق أن يطلب عملاؤها زيارة فرع عبر البوابة، فينشئ مسؤول النظام خدمة «طلب زيارة فرع» حول النموذج المنشور بالاسم نفسه، ويجعل الزيارة خطوة حضورية في فرع حلب. ثم يسجّل عمر حداد الطلبات الحضورية في مكتب الاستقبال، وتتابع سارة الخطيب أداء الخدمة وما تأخر من طلباتها.
\end{rcnote}

\section{صفحات هذا الفصل}
\label{sec:builder-services-pages}

تبدأ من صفحة «الخدمات» في قسم «الأساسية» من أدوات البناء، ومنها تفتح كل خدمة في صفحة لها خمسة تبويبات. وترتيب الصفحات هنا هو ترتيب البناء المقترح، فاتبعه في المرة الأولى، ثم عُد إلى أي صفحة عندما تحتاج إليها.

\begin{table}[H]
\centering
\caption{تقديم الخدمات خطوة بخطوة}
\label{tab:builder-pages-services}
\rcZebra
\begin{tabularx}{\linewidth}{R{4.2cm} Y}
\rcHeaderRow
\thh{الصفحة وقسمها} & \thh{ماذا تقدم هذه الصفحة؟}\\
الخدمات — الأساسية & تعرض الخدمات المعروضة على البوابة والنماذج الداخلية، ومنها تنشئ خدمة جديدة وتختار طريقة تقديمها.\\
نظرة عامة على الخدمة — صفحة الخدمة & تلخّص جاهزية الخدمة وأرقامها في آخر 30 يوماً والطلبات المتأخرة وسجل تغييراتها.\\
التعريف ودليل الخدمة — صفحة الخدمة & اسم الخدمة ووصفها وعنوانها على البوابة، والدليل الذي يقرؤه المتعامل قبل أن يبدأ.\\
التقديم والأهلية — صفحة الخدمة & طريقة التقديم وحالة الخدمة، ومن يحق له التقديم وبأي شروط، وحدود الطلبات ومعالجتها.\\
مراكز الخدمة — صفحة الخدمة & الأماكن التي تجري فيها الخطوات الحضورية، وأوقات دوامها، والموظفون الذين يستقبلون فيها.\\
المسار والنماذج — صفحة الخدمة & الخطوات التي يمر بها الطلب بالترتيب، والنماذج خارج المسار، ونماذج المتابعة.\\
الرسوم — صفحة الخدمة & بنود الرسوم ومتى تُستحق، والإعفاءات، وطريقة الدفع وما يحدث للطلب عند الاستحقاق.\\
إتاحة نموذج على البوابة — الخدمات & تحوّل نموذجاً داخلياً منشوراً إلى خدمة جديدة أو إلى خطوة في خدمة قائمة.\\
الخدمة في إعدادات النموذج — منشئ النماذج & تبيّن داخل النموذج الخدمة التي ينتمي إليها، وما تفرضه عليه، ومدى توافقهما.\\
مكتب الاستقبال — الخدمات & يسجّل فيه موظف المركز الطلب الحضوري ويوثّق المستندات المستلمة ويطبع إيصالاً برقم الطلب.\\
أداء الخدمات — الخدمات & يقيس عدد الطلبات ومدد إنجازها، ويبيّن المراحل المتأخرة ورضا المتعاملين.\\
\end{tabularx}
\end{table}
